\documentclass[11pt,a4paper]{article}

\usepackage[margin=2.2cm]{geometry}
\usepackage{amsmath,amssymb,amsthm}
\usepackage{booktabs}
\usepackage{xcolor}
\usepackage{hyperref}
\usepackage{parskip}
\usepackage{lmodern}

\hypersetup{colorlinks=true, linkcolor=black, urlcolor=blue, citecolor=blue}

\newtheorem{lemma}{Lemma}
\newtheorem{theorem}{Theorem}

\title{A Cross-Layer Mutual-Information Extension\\[2pt]
       via Matrix-Based R\'enyi Entropy}
\author{thesis-IV working notes (HyMeKo)}
\date{2026-04-26}

\newcommand{\trace}{\operatorname{tr}}
\newcommand{\Khat}{\hat{K}}

\begin{document}
\maketitle

\begin{abstract}
We formalise the cross-layer mutual-information regulariser proposed as
Path~F of the universality programme. Using the matrix-based R\'enyi-$\alpha$
entropy framework of Sanchez-Giraldo, Rao, and Principe, we show that
four extensions of the existing scalar-spectral-entropy regulariser
(temporal-conditional, cross-layer mutual information, joint view
entropy, and matrix-based information bottleneck) are all instances of
a single primitive: the joint of two trace-normalised activation Gram
matrices via Hadamard product. The $\alpha=2$ collision-entropy form
admits a closed-form, eigendecomposition-free, batch-quadratic estimator
and is differentiable through standard autograd. Three falsifiable
predictions for the phase-8 empirical test conclude.
\end{abstract}

\section{Motivation}

The nine entropy regularisers tested in phases 1--7b all use
\emph{marginal} spectral entropy at a single timestep. Specifically, for
layer weights $W$ with adjacency $A$ and normalised-Laplacian eigenvalues
$\{\lambda_i\}$,
\[
H_{\text{marg}}(A) = -\sum p_i \log p_i,
\qquad p_i = \lambda_i / \textstyle\sum_j \lambda_j .
\]
This captures \emph{how spread out} the weight spectrum is, but never:
\begin{itemize}
  \item the \textbf{joint} distribution across layers $H(W_i, W_j)$,
  \item the \textbf{conditional} distribution across time
        $H(W_t \mid W_{t-1})$,
  \item or the \textbf{mutual information} between any of those.
\end{itemize}
Phases~6 and~7 demonstrated that a single normalisation
(\texttt{scalar\_entropy\_normalized}, Path~B) and target-entropy
formulation (\texttt{entropy\_target}, Path~A) achieve universality on a
4-dataset stress matrix (spirals, circles, MNIST plain MLP, CapsMLP MNIST),
but neither addresses the structural hypothesis that \emph{layers may be
redundant with each other} and that decorrelating them might be a stronger
universality lever than tuning a single scalar target.

This note develops a single quantity --- the matrix-based R\'enyi-$\alpha$
mutual information between activation Gram matrices of two layers --- which
subsumes the four extension paths under one definition.

\section{Setup and notation}

Let $M : \mathbb{R}^d \to \mathbb{R}^c$ be a feedforward network with $L$
layers parameterised by weights $W_1, \dots, W_L$. Let $X \in
\mathbb{R}^{B \times d}$ be a batch of $B$ input examples. Define the
per-layer activation tensor
\[
a_l \;=\; \varphi_l\bigl( W_l\, \varphi_{l-1}(\cdots W_1 X^\top) \bigr)^\top
   \;\in\; \mathbb{R}^{B \times d_l},
\]
i.e.\ the post-non-linearity output of layer~$l$ evaluated on $X$.

Construct the \textbf{sample-side Gram matrix} of layer~$l$,
\[
K_l(X) = a_l a_l^\top \in \mathbb{R}^{B \times B},  \tag{1}
\]
normalised so trace is~1,
\[
\Khat_l = K_l / \trace(K_l), \qquad \trace(\Khat_l) = 1. \tag{2}
\]
$\Khat_l$ is symmetric positive semi-definite and trace-$1$; equivalently,
its eigenvalue distribution $\{\mu_i^{(l)}\}$ is a valid probability mass
function on $B$ outcomes.

\section{Matrix-based R\'enyi-$\alpha$ entropy}

For $\alpha > 0$, $\alpha \neq 1$, define
\begin{equation}
H_\alpha(\Khat) \;=\; \frac{1}{1-\alpha} \log \sum_i \mu_i^\alpha
            \;=\; \frac{1}{1-\alpha} \log \trace(\Khat^\alpha).  \tag{3}
\end{equation}

\paragraph{Special cases.}
$\alpha \to 1$ recovers the Shannon (von Neumann) entropy
$H_1(\Khat) = -\sum \mu_i \log \mu_i$.

$\alpha = 2$ gives the \textbf{collision entropy}
\[
H_2(\Khat) = -\log \sum_i \mu_i^2 = -\log \trace(\Khat^2).  \tag{4}
\]
Equation~(4) is the workhorse: it requires no eigendecomposition, scales
as $O(B^2)$ from $\trace(\Khat^2) = \sum_{i,j} \Khat_{ij}^2$, and is
automatically differentiable in any standard autograd framework.

\section{Joint distribution: Hadamard product}

For two trace-$1$ PSD matrices $\Khat_i, \Khat_j \in \mathbb{R}^{B
\times B}$, define the \textbf{joint matrix} as the entrywise (Hadamard)
product, normalised:
\begin{equation}
\Khat_{ij} \;=\; (\Khat_i \odot \Khat_j) /
                   \trace(\Khat_i \odot \Khat_j). \tag{5}
\end{equation}

\begin{lemma}[Schur product theorem]
$\Khat_i \odot \Khat_j$ is positive semi-definite.
\end{lemma}
\begin{proof}[Sketch]
The Hadamard product of two PSD matrices is PSD; this is the classical
Schur theorem. Both factors are PSD by construction (eq.~1).
\end{proof}

\begin{lemma}[Independence factorisation]
If $\Khat_i$ and $\Khat_j$ correspond to statistically independent
features over the batch, then $H_\alpha(\Khat_{ij}) = H_\alpha(\Khat_i) +
H_\alpha(\Khat_j)$.
\end{lemma}
\begin{proof}[Sketch]
Direct manipulation of $\trace(\Khat^\alpha)$ using
spectral-factor independence; full proof in Sanchez-Giraldo et al.\
(2014), Theorem~5.
\end{proof}

\section{Cross-layer mutual information}

Define
\begin{equation}
I_\alpha(\Khat_i ; \Khat_j) \;=\; H_\alpha(\Khat_i) + H_\alpha(\Khat_j)
                                  - H_\alpha(\Khat_{ij}). \tag{6}
\end{equation}

\begin{theorem}[Sanchez-Giraldo, Theorem~6]
For $\alpha \in (0,1) \cup (1, \infty)$, $I_\alpha(\Khat_i ; \Khat_j)
\geq 0$, with equality iff $\Khat_{ij} = \Khat_i \odot \Khat_j$ holds in
the independence-factorisation sense.
\end{theorem}
\begin{proof}[Sketch]
Apply the log-sum inequality to the spectrum of $\Khat_{ij}$ relative
to the product of marginal spectra. The Hadamard product satisfies
\[
\trace\bigl((\Khat_i \odot \Khat_j)^\alpha\bigr)
   \;\leq\; \trace(\Khat_i^\alpha)^{1/2} \cdot \trace(\Khat_j^\alpha)^{1/2}
\]
(Cauchy-Schwarz on traces); rearranging gives the bound.
\end{proof}

\paragraph{Properties.}
\begin{itemize}
  \item Symmetry: $I_\alpha(\Khat_i ; \Khat_j) = I_\alpha(\Khat_j ; \Khat_i)$
        --- immediate from Hadamard commutativity.
  \item Boundedness: $0 \leq I_\alpha(\Khat_i ; \Khat_j) \leq
        \min\{H_\alpha(\Khat_i), H_\alpha(\Khat_j)\}$.
\end{itemize}

\section{Conditional entropy}

Define
\begin{equation}
H_\alpha(\Khat_j \mid \Khat_i)
   \;=\; H_\alpha(\Khat_{ij}) - H_\alpha(\Khat_i). \tag{7}
\end{equation}

\paragraph{Key properties.}
\begin{itemize}
  \item $H_\alpha(\Khat_j \mid \Khat_i) \geq 0$.
  \item $H_\alpha(\Khat_j \mid \Khat_i) = H_\alpha(\Khat_j) -
        I_\alpha(\Khat_i ; \Khat_j)$ (chain rule).
  \item Chain rule recovers eq.~(6).
\end{itemize}

\section{Differentiability and computational complexity}

\paragraph{Differentiability.} With $\alpha = 2$,
\[
H_2(\Khat) = -\log \trace(\Khat^2),
\quad \trace(\Khat^2) = \sum_{i,j} \Khat_{ij}^2,
\quad \Khat = (a a^\top) / \trace(a a^\top).
\]
Each operation is a smooth function of the activations $a$, hence of the
weights through the chain rule. Standard autograd handles this without
manual derivatives.

For $\alpha \neq 2$ we need eigenvalues; \texttt{torch.linalg.eigvalsh}
provides gradients via the Bunch-Lehmberger formula on symmetric inputs.

\paragraph{Computational cost per regulariser-update step.}
\begin{center}
\begin{tabular}{@{}lr@{}}
\toprule
Step & Cost \\
\midrule
Compute $\Khat_l$ for one layer            & $O(B^2 d_l)$ \\
Compute $H_2(\Khat_l)$ per layer            & $O(B^2)$ \\
Compute $\Khat_{ij}$ per pair               & $O(B^2)$ \\
Compute $H_2(\Khat_{ij})$ per pair          & $O(B^2)$ \\
Number of pairs                             & $L(L{-}1)/2$ \\
\midrule
\textbf{Total at $\alpha = 2$}              & $\bm{O(L^2 B^2)}$ per step \\
\bottomrule
\end{tabular}
\end{center}

For our setups $B = 128$ and $L \in \{3, \dots, 21\}$ this is
${<}1$\,ms on GPU per regulariser update.

\section{The four extensions, unified}

\paragraph{Path~E --- temporal conditional entropy.}
For batch $X$ shared across consecutive regulariser updates with
layer-$l$ Gram $\Khat_l^{(t)}$ and $\Khat_l^{(t-1)}$,
\[
\mathcal{L}_E = \lambda \cdot H_\alpha\bigl(\Khat_l^{(t)} \,\big|\,
                                                \Khat_l^{(t-1)}\bigr).
\]
Penalises \emph{novel} layer information per step.

\paragraph{Path~F --- cross-layer mutual information.}
\[
\mathcal{L}_F = \lambda \cdot \sum_{i < j} I_\alpha(\Khat_i ; \Khat_j).
\]
Penalises redundancy between layers; encourages each layer to encode
information not already in its predecessors.

\paragraph{Path~G --- joint view entropy.}
For dataflow- and factor-view adjacencies $K_{\text{df}}, K_{\text{ft}}$,
\[
\mathcal{L}_G = \lambda \cdot H_\alpha
                \bigl(\Khat_{\text{df}} \odot \Khat_{\text{ft}}\bigr).
\]
Removes the user's choice of view.

\paragraph{Path~H --- information bottleneck.}
For input/target batches with Gram $\Khat_X, \Khat_Y$ and a hidden
representation $T = a_l(X)$ with Gram $\Khat_T$,
\[
\mathcal{L}_H = \lambda \cdot I_\alpha(\Khat_X ; \Khat_T)
                - \beta \cdot I_\alpha(\Khat_T ; \Khat_Y).
\]
The matrix-based information-bottleneck of Achille and Soatto.

\paragraph{Unification.} Paths E, F, G, H are all functions of the same
primitive $I_\alpha(\Khat_a ; \Khat_b)$, with differing arguments
(layer/time/view/data axes). Choosing one implementation gives all four
extensions for free.

\section{Caveats and design choices}

\begin{enumerate}
  \item \textbf{Data dependence.} Unlike the existing weight-only
        regularisers, $\Khat_l$ requires a forward pass. In our training
        loop this is free (one forward already happens for the task
        loss). \emph{However}, the regulariser becomes stochastic across
        mini-batches.
  \item \textbf{Mini-batch noise mitigation.} Sample a fixed held-out
        batch $X_{\text{reg}}$ once at training start and compute
        $\Khat_l(X_{\text{reg}})$. Removes mini-batch fluctuations.
  \item \textbf{Choice of $\alpha$.} $\alpha = 2$ has the closed-form
        $-\log \trace(\Khat^2)$ and is the recommended default for
        compute. Shannon ($\alpha\to1$) is the ``standard'' entropy
        reviewers expect; we default to $\alpha = 2$ and report
        Shannon as a sensitivity analysis.
  \item \textbf{Activation hook point.} Hook \textbf{post-linear,
        pre-activation} to keep the closest connection to the existing
        weight-side adjacency; the post-non-linear hook is also
        defensible and slightly cheaper.
\end{enumerate}

\section{Hypotheses for empirical testing (phase~8)}

If layer redundancy is the structural cause of the residual negatives
(circles, CapsMLP) under the original weight-side regulariser, we
predict:
\begin{enumerate}
  \item \textbf{Positive datasets} (spirals, MNIST plain MLP) should
        have non-negligible cross-layer MI; the MI penalty should
        preserve the positive sign --- with possibly smaller magnitude
        than scalar entropy (since MI penalises a different structural
        quantity).
  \item \textbf{Negative datasets} (circles, CapsMLP MNIST) should have
        \emph{higher} cross-layer MI than positives --- the failure mode
        is layers learning redundant information instead of a
        distributed code. If this prediction holds, decorrelating via MI
        penalty should \emph{neutralise} the negatives more cleanly than
        scalar-entropy normalisation does.
  \item \textbf{Variance reduction}: the dominant cross-cutting effect
        of the weight-side family was $\sigma$-reduction in val-acc
        across seeds. We predict the activation-side MI penalty also
        reduces $\sigma$, possibly \emph{more strongly} on architectures
        with redundant pathways (CapsMLP, ResMLP-20).
\end{enumerate}

A clear pass through these three predictions would justify presenting
the activation-side framework as the principled refinement of the
weight-side one in the manuscript.

\bibliographystyle{plain}
\begin{thebibliography}{9}
\bibitem{sanchezgiraldo2014}
Sanchez-Giraldo, L., Rao, M., and Principe, J.~C.
\newblock {\em Measures of Entropy from Data Using Infinitely Divisible
  Kernels.}
\newblock IEEE Trans.\ Inform.\ Theory 61(1), 2014.

\bibitem{yu2019}
Yu, S., Sanchez-Giraldo, L., and Principe, J.~C.
\newblock {\em Multivariate Extension of Matrix-Based R\'enyi's
  $\alpha$-Order Entropy Functional.}
\newblock IEEE Trans.\ PAMI 42(8), 2019.

\bibitem{wickstrom2020}
Wickstr\o m et al.
\newblock {\em Information Plane Analysis of Deep Neural Networks via
  Matrix-Based R\'enyi's Entropy and Tensor Kernels.}
\newblock JMLR 21, 2020.

\bibitem{achillesoatto2018}
Achille, A. and Soatto, S.
\newblock {\em Information Dropout: Learning Optimal Representations
  Through Noisy Computation.}
\newblock IEEE Trans.\ PAMI 40(12), 2018.

\bibitem{tishby2015}
Tishby, N. and Zaslavsky, N.
\newblock {\em Deep Learning and the Information Bottleneck Principle.}
\newblock ITW, 2015.
\end{thebibliography}

\end{document}
